\chapter{Algorithmes d'appariement et spreads implicites}\label{ch:m1:matching-algorithms-and-implied-spreads}

Deux traders affichent le même cours acheteur pour le même contrat à terme. Un ordre de vente arrive,
de taille inférieure à leur taille cumulée. Sur le contrat à terme sur indice d'actions, le trader
qui a affiché le premier son cours acheteur est exécuté entièrement avant que l'autre ne reçoive la moindre taille. Sur
le contrat à terme sur taux d'intérêt à court terme voisin, sur la même bourse, les deux
se partagent la vente au prorata des tailles qu'ils ont affichées, et celui qui
est arrivé le premier n'a aucun avantage. La première règle récompense la rapidité~; la
seconde récompense la taille, et pousse chacun à afficher davantage qu'il ne le souhaite. La
règle d'allocation tient en quelques lignes de la documentation de la bourse, et elle détermine
du type de firme qui peut faire office de \omterm{def:m1:what-a-trading-firm-does:market-maker}{teneur de marché} sur le produit. Un second mécanisme,
invisible à l'écran, relie les carnets de différentes échéances~: le meilleur cours acheteur
d'un \omterm{def:m1:matching-algorithms-and-implied-spreads:calendar}{calendar spread} est souvent un ordre que personne n'a saisi.

\section{Priorité prix-temps}

\begin{definition}[Priorité prix-temps]\label{def:m1:matching-algorithms-and-implied-spreads:fifo}
Selon la \emph{priorité prix-temps}\index{priorité prix-temps} (premier entré, premier
sorti, FIFO), un ordre entrant s'exécute contre les ordres au repos au meilleur
prix, dans l'ordre de leurs horodatages de priorité, du plus ancien au plus récent, chacun étant
exécuté entièrement avant que le suivant ne reçoive quoi que ce soit.
\end{definition}

L'horodatage d'un ordre est fixé à son entrée dans le carnet. Sur la plupart des bourses, il est
renouvelé, et la place dans la file perdue, lorsque le prix de l'ordre est modifié
ou que sa quantité est augmentée~; la réduction de la quantité conserve la place. Sous FIFO,
un ordre au repos possède deux attributs qui comptent~: son prix et sa position
dans la file. Lorsque le tick est grand (\cref{fig:m1:futures-contracts-and-exchanges:tickbp}),
le prix change rarement, les files sont longues, et la position fait tout le jeu~:
être proche de la tête, c'est être exécuté par de petits ordres non informés~;
être en queue, c'est n'être exécuté que lorsqu'un gros ordre balaie le niveau,
ce qui intervient au moment où le prix est sur le point de jouer contre l'exécution. L'économie
des files d'attente est développée dans le livre~10.

\section{Allocation au prorata et variantes}

\begin{definition}[Allocation au prorata]\label{def:m1:matching-algorithms-and-implied-spreads:prorata}
Selon l'\emph{allocation au prorata}\index{allocation au prorata}, un ordre entrant
est partagé entre tous les ordres au repos au meilleur prix, proportionnellement à leurs
tailles. Les allocations sont arrondies au lot entier inférieur, les allocations inférieures à un
minimum propre au produit peuvent être supprimées, et ce que l'arrondi laisse en reste est
réparti selon une règle secondaire, souvent la priorité temporelle.
\end{definition}

\begin{definition}[Allocation au premier ordre et teneur de marché principal]\label{def:m1:matching-algorithms-and-implied-spreads:top}
Une \emph{allocation au premier ordre}\index{allocation au premier ordre} accorde à l'ordre qui a
le premier amélioré le marché, en établissant un nouveau meilleur prix, une allocation prioritaire
avant l'application de toute autre règle. Un \emph{teneur de marché
principal}\index{teneur de marché principal} (lead market maker) est une firme qui, en échange
d'obligations de cotation, reçoit un pourcentage fixé de chaque ordre entrant au meilleur
prix avant que le reste ne soit alloué.
\end{definition}

Les algorithmes réels sont des séquences de telles étapes. Une famille en usage~: l'ordre
au premier rang, jusqu'à un pourcentage de l'agresseur~; puis les \omterm{def:m1:what-a-trading-firm-does:market-maker}{teneurs de marché} principaux~;
puis un pourcentage de ce qui reste, par le temps~; puis le reste au prorata avec une
allocation minimale~; puis les restes par le temps.

\begin{dated}{2026-09}{Quelle règle pour quel produit}\label{dat:m1:matching-algorithms-and-implied-spreads:which}
\textbf{CME Group} documente neuf algorithmes, parmi lesquels FIFO, FIFO avec
\omterm{def:m1:matching-algorithms-and-implied-spreads:top}{teneur de marché principal}, Pro Rata, Allocation (premier ordre, puis prorata avec un
minimum de deux lots, puis FIFO), Threshold Pro Rata (avec et sans \omterm{def:m1:matching-algorithms-and-implied-spreads:top}{teneur de marché
principal}) et Configurable, une répartition de FIFO et de prorata selon des
pourcentages fixés. Une enquête d'un fournisseur de données sur les définitions de la bourse attribue
FIFO aux contrats au comptant sur indice d'actions, aux obligations du Trésor à dix ans et au pétrole brut, soit environ 70~\%
du volume~; Configurable aux obligations à deux ans, aux fed funds et aux céréales,
soit environ 13~\%~; Allocation aux contrats au comptant SOFR à trois mois, soit environ 10~\%~; et
Threshold Pro Rata avec \omterm{def:m1:matching-algorithms-and-implied-spreads:top}{teneur de marché principal} aux options sur l'obligation à dix ans.
\textbf{Eurex} prend en charge l'allocation par le temps, au prorata et par le temps puis au prorata, cette dernière
donnant aux ordres plus anciens davantage que leur part proportionnelle~; le 30 mars 2026, elle a
fait passer ses contrats à terme sur le marché monétaire, dont l'Euribor à trois mois, du
temps à l'allocation par le temps puis au prorata.
\end{dated}

La logique est économique. Les produits dont le prix se déplace de nombreux ticks par jour conservent
FIFO~: le tick est petit, améliorer le prix est le moyen d'obtenir la priorité, et les
files sont courtes. Les contrats à terme sur taux d'intérêt à court terme se déplacent d'un ou deux ticks par jour~:
sous FIFO, la file à chaque prix serait longue de plusieurs jours, et personne d'autre
que le premier arrivé ne prendrait la peine de coter. Le partage des exécutions au prorata
maintient de nombreuses firmes prêtes à afficher de la taille à un prix qui ne bouge pas.

\begin{omfigure}
\begin{tikzpicture}
\begin{axis}[width=11cm,height=5.2cm,ybar,bar width=9pt,
  ylabel={lots reçus},
  symbolic x coords={A,B,C,D},xtick=data,
  xticklabels={{A~: 10, premier}, {B~: 200}, {C~: 40, TMP}, {D~: 250, dernier}},
  tick label style={font=\small},label style={font=\small},xticklabel style={font=\footnotesize\strut},
  ymin=0,ymax=105,ymajorgrids,grid style={gray!25},enlarge x limits=0.15,
  nodes near coords,nodes near coords style={font=\scriptsize},
  legend columns=1,legend cell align=left,legend style={font=\footnotesize,at={(0.5,-0.2)},anchor=north,draw=none,fill=none,/tikz/every even column/.append style={column sep=8pt}},
  table/col sep=comma]
\addplot[fill=omDef!70,draw=omDef] table[x=order,y=fifo]{figdata/markets-1/19-matching-algorithms-and-implied-spreads/alloc.csv};
\addplot[fill=omThm!60,draw=omThm] table[x=order,y=pro_rata]{figdata/markets-1/19-matching-algorithms-and-implied-spreads/alloc.csv};
\addplot[fill=omExo!50,draw=omExo] table[x=order,y=split]{figdata/markets-1/19-matching-algorithms-and-implied-spreads/alloc.csv};
\legend{FIFO,prorata (minimum 2),répartition~: premier 10~\% / TMP 10~\% / 40~\% par le temps / reste au prorata}
\end{axis}
\end{tikzpicture}

\omcaption{Cent lots arrivent à un niveau comportant quatre ordres, classés en
priorité temporelle avec leurs tailles. Le même agresseur, le même carnet, trois
règles, trois ensembles de gagnants différents. Données~: la
construction du chapitre.}\label{fig:m1:matching-algorithms-and-implied-spreads:alloc}
\end{omfigure}

\begin{proposition}[Le prorata récompense la taille affichée]\label{prop:m1:matching-algorithms-and-implied-spreads:oversize}
Sous prorata sans minimum, un ordre de taille $s$ dans un niveau dont les autres
ordres totalisent $Q$ reçoit $\lfloor as/(Q+s)\rfloor$ d'un agresseur de taille $a
\le Q+s$, croissant en $s$ et indépendant du moment d'arrivée. Un trader qui
souhaite obtenir $w$ lots d'un agresseur typique $a$ doit afficher environ $s = wQ/(a-w)$,
et il s'expose à recevoir jusqu'à $s$ lorsqu'un agresseur plus grand arrive.
\end{proposition}

\begin{example}[Quatre-vingt-sept pour en obtenir huit]\label{ex:m1:matching-algorithms-and-implied-spreads:87}
Un niveau contient 1\,000 lots et les agresseurs typiques portent sur 100 lots. Pour recevoir 8
le trader affiche $8 \times 1\,000/92 = 87$ lots (86 en recevrait 7). Si un
ordre de 1\,000 lots arrive à la place, le trader en reçoit 80~: dix fois
son intention,
au moment où le prix risque le plus de traverser le niveau.
La taille affichée dans un carnet au prorata surestime donc la quantité que
chacun souhaite négocier, et une grande part disparaît lorsqu'un gros ordre arrive.
\end{example}

\begin{omfigure}
\begin{tikzpicture}
\begin{axis}[width=10cm,height=4.6cm,
  xlabel={taille affichée par le nouvel arrivant (lots)},ylabel={lots reçus},
  tick label style={font=\small},label style={font=\small},
  xmin=0,xmax=400,ymin=0,ymax=32,grid=major,grid style={gray!25},table/col sep=comma]
\addplot[omDef,very thick,const plot] table[x=shown,y=filled]{figdata/markets-1/19-matching-algorithms-and-implied-spreads/oversize.csv};
\end{axis}
\end{tikzpicture}

\omcaption{Prorata avec une allocation minimale de deux lots~: ce que reçoit un nouvel arrivant
en queue d'un niveau de 1\,000 lots d'un agresseur de 100 lots, en fonction
de la taille qu'il affiche. En dessous de 21 lots, il ne reçoit rien. Données~:
la \cref{prop:m1:matching-algorithms-and-implied-spreads:oversize}.}\label{fig:m1:matching-algorithms-and-implied-spreads:oversize}
\end{omfigure}

\begin{omfigure}
\begin{tikzpicture}
\begin{axis}[width=10.5cm,height=4.8cm,
  xlabel={lots devant dans la file},ylabel={exécution espérée (lots)},
  tick label style={font=\small},label style={font=\small},
  xmin=0,xmax=490,ymin=0,ymax=10.5,grid=major,grid style={gray!25},
  legend columns=2,legend style={font=\footnotesize,at={(0.5,-0.32)},anchor=north,draw=none,fill=none,/tikz/every even column/.append style={column sep=8pt}},
  table/col sep=comma]
\addplot[omDef,very thick] table[x=ahead,y=fifo]{figdata/markets-1/19-matching-algorithms-and-implied-spreads/position.csv};
\addplot[omThm,very thick,dashed] table[x=ahead,y=pro_rata]{figdata/markets-1/19-matching-algorithms-and-implied-spreads/position.csv};
\legend{FIFO,prorata}
\end{axis}
\end{tikzpicture}

\omcaption{Exécution espérée d'un ordre de 10 lots dans un niveau de 500 lots lorsqu'un
agresseur de taille aléatoire (exponentielle, moyenne 120) arrive. Sous FIFO, le
début de la file vaut quatre fois la part au prorata et la fin presque
rien~; les courbes se croisent avec 170 lots devant. Données~:
la simulation du tutoriel.}\label{fig:m1:matching-algorithms-and-implied-spreads:position}
\end{omfigure}

\section{Calendar spreads et prix implicites}

\begin{definition}[Calendar spread]\label{def:m1:matching-algorithms-and-implied-spreads:calendar}
Un \emph{calendar spread}\index{calendar spread} est un instrument coté en bourse
dont l'achat consiste à acheter simultanément une échéance et à vendre une autre
du même produit, à une différence de prix cotée. Il possède
son propre carnet d'ordres, son propre tick, et se traite comme un seul ordre, sans risque
de n'être exécuté que sur une seule jambe.
\end{definition}

La convention retenue ici~: le spread est égal à l'échéance la plus proche moins
la plus lointaine, de sorte qu'acheter le spread achète l'échéance la plus proche et vend la plus lointaine. Faire le roll d'une position longue
vers l'échéance suivante revient à vendre le spread.

\begin{definition}[Implied-in et implied-out]\label{def:m1:matching-algorithms-and-implied-spreads:implied}
Un ordre \emph{implied-in}\index{implied-in} est un ordre de spread que le
moteur d'appariement déduit d'ordres au repos dans les carnets outright~: un cours acheteur
sur l'échéance la plus proche et un cours vendeur sur l'échéance la plus lointaine impliquent un cours acheteur sur le spread
au différentiel de leurs prix, pour le plus petit de leurs deux volumes. Un
\emph{implied-out}\index{implied-out} est un ordre outright dérivé d'un
ordre sur spread et d'un ordre outright sur l'autre jambe. Les ordres implicites sont
réels~: ils peuvent être touchés, et le moteur exécute alors toutes les jambes d'un seul coup.
Les ordres implicites ne s'échangent pas entre eux.
\end{definition}

\begin{proposition}[Prix implicites]\label{prop:m1:matching-algorithms-and-implied-spreads:implied}
Soient $b_F, a_F$ et $b_B, a_B$ les meilleurs cours acheteur et vendeur des
échéances la plus proche et la plus lointaine, et $b_S, a_S$ ceux saisis directement dans le spread~:
\begin{align*}
\text{implied-in~:}\quad & b_S^{\text{imp}} = b_F - a_B, & a_S^{\text{imp}} &= a_F - b_B,\\
\text{implied-out, échéance la plus proche~:}\quad & b_F^{\text{imp}} = b_S + b_B, & a_F^{\text{imp}} &= a_S + a_B,\\
\text{implied-out, échéance la plus lointaine~:}\quad & b_B^{\text{imp}} = b_F - a_S, & a_B^{\text{imp}} &= a_F - b_S.
\end{align*}
La taille de chaque ordre implicite est la plus petite des tailles de ses deux composantes.
\end{proposition}
\begin{proof}
Un vendeur du spread vend l'échéance la plus proche et achète la plus lointaine. Le moteur peut
servir ce vendeur avec l'acheteur sur l'échéance la plus proche et le vendeur sur la plus lointaine, qui
paient ensemble $b_F$ et facturent $a_B$. Les cinq autres cas relèvent du même raisonnement appliqué
à chaque jambe.
\end{proof}

\begin{omfigure}
\begin{tikzpicture}[font=\small,
  bk/.style={draw=omDef,rounded corners=2pt,minimum width=3.3cm,minimum height=1.5cm,align=center,fill=omDef!4},
  flow/.style={-{Stealth[length=2mm]},thick}]
\node[bk] (f) at (0,1.2) {\textbf{Échéance la plus proche}\\bid 10\,050 $\times$ 30\\ask 10\,052 $\times$ 25};
\node[bk] (b) at (0,-1.2) {\textbf{Échéance la plus lointaine}\\bid 10\,020 $\times$ 40\\ask 10\,023 $\times$ 10};
\node[bk,draw=omThm,fill=omThm!5,minimum width=5.4cm,minimum height=2.3cm] (s) at (8.0,0) {\textbf{Spread (proche $-$ lointaine)}\\direct~: bid 28 $\times$ 5, ask 32 $\times$ 7\\implicite~: bid 27 $\times$ 10, ask 32 $\times$ 25\\[2pt]\emph{affiché}~: bid 28 $\times$ 5, ask 32 $\times$ 32};
\draw[flow,omThm] (f.east) -- node[above,sloped,font=\footnotesize]{bid $-$} (s.west|-0,0.35);
\draw[flow,omThm] (b.east) -- node[below,sloped,font=\footnotesize]{$-$ ask} (s.west|-0,-0.35);
\end{tikzpicture}

\omcaption{Ordres \omterm{def:m1:matching-algorithms-and-implied-spreads:implied}{implied-in}. Le bid implicite du spread est le bid de l'échéance proche
moins l'ask de l'échéance lointaine, $10\,050 - 10\,023 = 27$, pour $\min(30, 10) = 10$ lots~; son
ask implicite est $10\,052 - 10\,020 = 32$ pour 25. Le carnet vu par un trader est le
meilleur entre le direct et l'implicite de chaque côté, les tailles s'additionnant à prix égal
prix.}\label{fig:m1:matching-algorithms-and-implied-spreads:implied}
\end{omfigure}

L'appariement implicite réunit la liquidité de toutes les échéances~: un hedger qui fait un
roll d'une position rencontre, par le spread, tous ceux qui cotent l'une ou l'autre échéance. Cela
signifie aussi qu'un ordre saisi dans un carnet peut s'exécuter à cause d'un événement survenu dans
un autre, que la position effective d'un ordre dans la file dépend de la règle
du moteur pour classer les ordres implicites face aux ordres directs au même prix, et qu'un
simulateur qui ignore les implicites sous-estime la liquidité du spread et
fausse les exécutions sur les outrights.

\section{Ce que montre le flux}

\begin{definition}[Carnet par prix et carnet par ordre]\label{def:m1:matching-algorithms-and-implied-spreads:mbo}
Un flux \emph{market-by-price}\index{market-by-price} (carnet par prix) publie, pour chaque
niveau de prix, la quantité totale et le nombre d'ordres. Un flux
\emph{market-by-order}\index{market-by-order} (carnet par ordre) publie chaque ordre
individuellement, avec un identifiant et sa priorité, ainsi que chaque modification le concernant.
\end{definition}

Avec un carnet par ordre, une firme connaît exactement sa position dans une file FIFO
et voit qui la précède en taille~; avec un carnet par prix, elle doit estimer sa
position à partir de la suite des mises à jour, en devinant si chaque baisse de
quantité était une annulation (devant ou derrière~?) ou une transaction. Les quantités implicites
sont généralement publiées séparément du carnet direct, ou pas du tout
au-delà du meilleur niveau. La lecture de ces flux fait l'objet du
\cref{ch:m1:reading-market-data}~; ce qui importe ici est que la règle d'allocation
détermine quel flux vaut la peine d'être payé.

\section{Tutoriel~: trois règles d'allocation}\label{tut:m1:matching-algorithms-and-implied-spreads:alloc}

\begin{tutorial}
\textbf{Objectif.} Implémenter la FIFO, le prorata avec un minimum et un partage
paramétrable~; comparer leurs exécutions~; calculer les prix implicites du spread. \textbf{État
final~:} les trois figures de données de ce chapitre.
\begin{tutsteps}
\item \textbf{FIFO}, et une fonction auxiliaire qui ne sur-remplit jamais un ordre.
\omcode{code/firm/match/firm_match.py}{18}{30}{Priorité temporelle~: parcourir la file.}\label{lst:m1:matching-algorithms-and-implied-spreads:fifo}
\item \textbf{Prorata.} Arithmétique entière uniquement~: parties entières, un minimum, et
restes répartis selon le temps.
\omcode{code/firm/match/firm_match.py}{38}{51}{Allocation au prorata sur les quantités ouvertes.}\label{lst:m1:matching-algorithms-and-implied-spreads:prorata}
\item \textbf{\omterm{def:m1:matching-algorithms-and-implied-spreads:implied}{Implied-in}.}
\omcode{code/firm/match/firm_match.py}{88}{93}{Cours de spread implicite issu de deux cotations outright.}\label{lst:m1:matching-algorithms-and-implied-spreads:implied}
\item \textbf{Tests de propriétés.} Sur trois cents carnets aléatoires, chaque algorithme
doit allouer exactement $\min(\text{agresseur}, \text{au repos})$ lots et jamais plus que
la taille d'un ordre.
\end{tutsteps}
\textbf{Pour aller plus loin.} Implémenter le prorata temporel~: pondérer chaque ordre par
sa taille multipliée par une fonction décroissante de son rang dans le temps, et trouver les poids
pour lesquels le premier ordre de la \cref{fig:m1:matching-algorithms-and-implied-spreads:alloc}
reçoit la même part que sous la règle de partage.
\end{tutorial}

\section{Construction~: allocation et implicites}\label{bld:m1:matching-algorithms-and-implied-spreads:match}

\begin{build}
\bfield{Rôle} Le moteur d'appariement de la bourse simulée de la firme miniature
(livre~10) délègue à ce module une seule décision~: étant donnés un niveau et
un agresseur, qui obtient quoi. Le backtester utilise le même code, de sorte que les
exécutions simulées respectent la règle du produit simulé.
\bfield{Interface} \texttt{Resting(oid, qty, lmm, top)} en priorité temporelle~;
\texttt{fifo(book, qty)}, \texttt{pro\_rata(book, qty, min\_alloc)},
\texttt{configurable(book, qty, top\_pct, lmm\_pct, fifo\_pct, min\_alloc)},
chacune renvoyant les exécutions par ordre~; \texttt{Quote}~; \texttt{implied\_in(front,
back)}, \texttt{implied\_out\_front(spread, back)}, \texttt{best\_of(direct,
implied)}.
\bfield{Règles} Entiers uniquement. Conservation~: les exécutions somment à $\min(\text{quantité},
\text{au repos})$. Aucun ordre n'est sur-rempli. Une quantité implicite ne s'apparie jamais
à une quantité implicite~: \texttt{best\_of} combine une cotation directe et une cotation implicite
et n'est jamais appliqué à deux cotations implicites.
\bfield{Tests d'acceptation} \texttt{code/firm/match/tests/}~: les allocations travaillées du
chapitre~; la règle de partage~; la conservation sur des carnets aléatoires~; les prix implicites
avec un côté vide.
\bfield{Pour aller plus loin} Un papillon (trois échéances) et la deuxième génération
d'implicites qu'il crée~; décider, comme doit le faire une bourse, où s'arrêter.
\end{build}

\begin{omsources}
\item CME Group Client Systems Wiki, \emph{Supported Matching Algorithms},
\emph{CME Globex Matching Algorithms} et \emph{Implied Orders}.
\item Databento, «~CME matching algorithms explained~» (étude des
algorithmes par produit et par part de volume).
\item Eurex, \emph{Matching principles}, et la circulaire sur les dérivés du
marché monétaire (schéma d'allocation à partir du 30 mars 2026).
\end{omsources}

\section{Exercices}

\begin{exercise}[$\star$]\label{exo:m1:matching-algorithms-and-implied-spreads:1}
Le niveau de la \cref{fig:m1:matching-algorithms-and-implied-spreads:alloc} (A 10,
B 200, C 40, D 250, dans l'ordre temporel) reçoit un agresseur de 230 lots sous
FIFO. Donner les exécutions.
\end{exercise}

\begin{exercise}[$\star$]\label{exo:m1:matching-algorithms-and-implied-spreads:2}
Même niveau, agresseur de 60 lots, prorata sans minimum et restes répartis selon
le temps. Donner les parts exactes, les parties entières, le reste et les exécutions.
\end{exercise}

\begin{exercise}[$\star$]\label{exo:m1:matching-algorithms-and-implied-spreads:3}
Échéance la plus proche~: cours acheteur 10\,050 pour 30, cours vendeur 10\,052 pour 25~; échéance
suivante~: cours acheteur 10\,020 pour 40, cours vendeur 10\,023 pour 10. Donner le cours acheteur et le
cours vendeur \omterm{def:m1:matching-algorithms-and-implied-spreads:implied}{implied-in} du spread, avec leurs tailles.
\end{exercise}

\begin{exercise}[$\star\star$]\label{exo:m1:matching-algorithms-and-implied-spreads:4}
Sous FIFO, lesquelles de ces actions sur un ordre au repos font perdre sa place dans la
file sur la plupart des bourses~: réduire sa taille~; augmenter sa taille~; modifier son
prix~; l'annuler puis le ressaisir~? Comment concevoir un gestionnaire d'ordres
qui doit \emph{augmenter} sa taille à un niveau sans renoncer à
la priorité qu'il possède~?
\end{exercise}

\begin{exercise}[$\star\star$]\label{exo:m1:matching-algorithms-and-implied-spreads:5}
Vérifier l'\cref{ex:m1:matching-algorithms-and-implied-spreads:87}. Supposer ensuite que
chaque participant raisonne de la même façon et multiplie sa taille par dix. Que se passe-t-il
pour les exécutions de chacun, pour la profondeur affichée et pour l'information
contenue dans la profondeur affichée~?
\end{exercise}

\begin{exercise}[$\star\star$]\label{exo:m1:matching-algorithms-and-implied-spreads:6}
Un flux \omterm{def:m1:matching-algorithms-and-implied-spreads:mbo}{market-by-price} montre le meilleur cours acheteur passant de 480 lots en 12 ordres à
465 lots en 11 ordres, sans message de transaction. Un ordre de 10 lots que l'on a placé
se trouve quelque part dans ce niveau, selon la FIFO. Que s'est-il passé, et que peut-on dire
et que ne peut-on pas dire de la position de cet ordre dans la file~? Qu'apporterait le \omterm{def:m1:matching-algorithms-and-implied-spreads:mbo}{market-by-order}~?
\end{exercise}

\begin{exercise}[$\star\star\star$]\label{exo:m1:matching-algorithms-and-implied-spreads:7}
\emph{Programmation.} À partir de \texttt{position.csv}~: indiquer l'exécution espérée en tête
de file selon la FIFO, l'exécution espérée au prorata, leur rapport, et le
nombre de lots devant soi au-delà duquel le prorata aurait été préférable.
\end{exercise}

\begin{exercise}[$\star\star\star$]\label{exo:m1:matching-algorithms-and-implied-spreads:8}
\emph{Repérer le défaut.} «~Le produit est au prorata avec un minimum de deux lots. Je
coterai 5 lots de chaque côté au meilleur prix, dans des niveaux d'environ 2\,000
lots, et récupérer ma juste part du flux, soit environ 0.25~\% de chaque transaction.~»
Les preneurs de liquidité typiques ont des ordres de 100 lots. Combien ce trader va-t-il récupérer, et pourquoi~?
\end{exercise}

\section{Problème~: le carnet implicite}

\begin{problem}[{Problème du week-end --- deux échéances, un spread, un moteur}]\label{pb:m1:matching-algorithms-and-implied-spreads:1}
Prix exprimés en ticks. Échéance la plus proche (Z)~: cours acheteur 10\,050 $\times$ 30, cours vendeur 10\,052
$\times$ 25. Échéance éloignée (H)~: cours acheteur 10\,020 $\times$ 40, cours vendeur 10\,023 $\times$ 10.
Ordres saisis directement dans le spread Z--H~: cours acheteur 28 $\times$ 5, cours vendeur 32 $\times$
7. Le spread est l'échéance la plus proche moins l'échéance éloignée.

\medskip\noindent\textbf{Partie I --- \omterm{def:m1:matching-algorithms-and-implied-spreads:implied}{Implied-in}.}
\begin{enumerate}
\item Donner le cours acheteur \omterm{def:m1:matching-algorithms-and-implied-spreads:implied}{implied-in} du spread, avec sa taille.
\item Donner le cours vendeur \omterm{def:m1:matching-algorithms-and-implied-spreads:implied}{implied-in}, avec sa taille.
\item Donner le carnet du spread que voit un trader au meilleur niveau.
\item Pourquoi le cours acheteur \omterm{def:m1:matching-algorithms-and-implied-spreads:implied}{implied-in} (27) ne peut-il pas rester au-dessus du cours
vendeur direct (32) pendant longtemps~? Que se passerait-il si c'était le cas~?
\item Expliquer pourquoi des ordres implicites ne doivent pas s'apparier avec d'autres ordres implicites.
\end{enumerate}

\medskip\noindent\textbf{Partie II --- Un roll.}
Un fonds vend 12 spreads au marché (il fait un roll de sa position longue vers l'échéance suivante).
\begin{enumerate}[resume]
\item Contre quels ordres le fonds s'exécute-t-il, et à quels prix~?
\item Donner son prix moyen d'exécution.
\item Quelles transactions sont publiées dans les carnets outright, et à quels prix~?
\item Donner les trois carnets au meilleur niveau après la transaction.
\item Le fonds aurait pu vendre 12 Z et acheter 12 H lui-même. Donner les prix qu'il
aurait obtenus à partir des carnets outright affichés, et l'écart.
\end{enumerate}

\medskip\noindent\textbf{Partie III --- \omterm{def:m1:matching-algorithms-and-implied-spreads:implied}{Implied-out} et un nouvel ordre.}
Revenir aux carnets initiaux.
\begin{enumerate}[resume]
\item Donner le cours acheteur et le cours vendeur \omterm{def:m1:matching-algorithms-and-implied-spreads:implied}{implied-out} sur l'échéance la plus proche, à partir des ordres directs
sur le spread et de l'échéance éloignée.
\item Améliorent-ils le marché affiché de l'échéance la plus proche~?
\item Un nouvel ordre de vente de 20 lots à 10\,021 entre sur l'échéance éloignée. Donner le nouveau
cours acheteur \omterm{def:m1:matching-algorithms-and-implied-spreads:implied}{implied-in} du spread.
\item Le nouvel ordre s'exécute-t-il~? Que faudrait-il pour cela~?
\item Un \omterm{def:m1:what-a-trading-firm-does:market-maker}{teneur de marché} affiche un cours acheteur de 50 à 28 dans le spread. Après l'ordre
de la question 13, où se situe-t-il par rapport au cours acheteur \omterm{def:m1:matching-algorithms-and-implied-spreads:implied}{implied-in}, et
qu'est-ce que cela lui apprend~?
\end{enumerate}

\medskip\noindent\textbf{Partie IV --- Jugement.}
\begin{enumerate}[resume]
\item Votre backtest d'une stratégie sur l'échéance la plus proche ne tient pas compte du carnet du spread. Citer
deux façons dont ses exécutions sont erronées.
\item Les carnets outright fonctionnent en FIFO. Un ordre \omterm{def:m1:matching-algorithms-and-implied-spreads:implied}{implied-out} et un ordre direct sont
affichés au même prix. Pourquoi la règle de classement du moteur entre les deux compte-t-elle
pour un \omterm{def:m1:what-a-trading-firm-does:market-maker}{teneur de marché}~?
\item Pourquoi les bourses ne publient-elles les quantités implicites qu'au meilleur niveau ou
aux deux premiers~?
\item Énoncer le \emph{résultat nommé}~: le meilleur cours acheteur \omterm{def:m1:matching-algorithms-and-implied-spreads:implied}{implied-in} du spread et
le meilleur cours acheteur affiché, avec leurs tailles.
\item En une phrase~: que fait l'appariement implicite pour un hedger~?
\end{enumerate}
\end{problem}

\section{Questions d'entretien}

\begin{interviewq}[$\star$ \iqroles{trader, développeur, chercheur}]\label{iq:m1:matching-algorithms-and-implied-spreads:1}
Expliquer l'appariement FIFO et l'appariement au prorata. Quels produits utilisent l'un ou l'autre, et pourquoi~?
\end{interviewq}

\begin{interviewq}[$\star$ \iqroles{développeur, trader}]\label{iq:m1:matching-algorithms-and-implied-spreads:2}
Qu'est-ce qu'un ordre implicite~?
\end{interviewq}

\begin{interviewq}[$\star\star$ \iqroles{trader, chercheur}]\label{iq:m1:matching-algorithms-and-implied-spreads:3}
Comment le comportement d'un \omterm{def:m1:what-a-trading-firm-does:market-maker}{teneur de marché} diffère-t-il entre un produit FIFO et un
produit au prorata~?
\end{interviewq}

\begin{interviewq}[$\star\star$ \iqroles{développeur}]\label{iq:m1:matching-algorithms-and-implied-spreads:4}
Écrire l'\omterm{def:m1:matching-algorithms-and-implied-spreads:prorata}{allocation au prorata} avec une taille de lot minimale. Quels sont les cas limites~?
\end{interviewq}

\begin{interviewq}[$\star\star$ \iqroles{chercheur, ingénieur ML}]\label{iq:m1:matching-algorithms-and-implied-spreads:5}
Vous ne disposez que de données \omterm{def:m1:matching-algorithms-and-implied-spreads:mbo}{market-by-price}. Comment estimez-vous votre position dans la file
dans un carnet FIFO~?
\end{interviewq}

\begin{interviewq}[$\star\star\star$ \iqroles{développeur, chercheur}]\label{iq:m1:matching-algorithms-and-implied-spreads:6}
Vous construisez un simulateur de bourse pour les \omterm{def:m1:matching-algorithms-and-implied-spreads:calendar}{calendar spreads}. Qu'est-ce que
l'appariement implicite vous oblige à bien traiter~?
\end{interviewq}
